\subsection{Typický scénář použití}
\label{subsec:kelvin-workflow}

Cyklus jedné úlohy v systému Kelvin zahrnuje několik na sebe navazujících kroků, od vytvoření zadání vyučujícím až po závěrečné manuální hodnocení odevzdaného řešení:

Vyučující nejprve vytvoří úlohu, kde zadá její textový popis, definuje testovací vstupy, očekávané výstupy a nastaví termín odevzdání.
Úloha je pak zpřístupněna přihlášeným studentům okamžitě nebo v předem stanoveném čase.

Student si zadání přečte, během přiděleného času vypracuje řešení ve svém vývojovém prostředí a výsledný zdrojový kód odevzdá prostřednictvím webového rozhraní systému.
Bezprostředně po odevzdání systém spustí automatizované testy.
Systém zkompiluje zdrojový kód a spustí jej proti připraveným testovacím vstupům, následné výsledky jsou porovnány s očekávanými výstupy.
Student okamžitě vidí výsledky testů které prošly i ty, které selhaly a následně může na základě zpětné vazby řešení opravit a odevzdat znovu.

Po uzavření termínu odevzdání přistoupí vyučující k manuálnímu hodnocení.
Prohlíží zdrojový kód odevzdaného řešení přímo v prostředí systému, přidává podle svého hodnocení komentáře ke konkrétním řádkům nebo blokům kódu a na základě celkového posouzení udělí bodové hodnocení.
Student hodnocení a komentáře obdrží a může na ně reagovat, čímž může vznikat diskuse nad zdrojovým kódem a jeho kvalitou.

\endinput
